We call $Z_g$ to be the surfaces obtained. A special case is $g=0$, where we define $Z_0$ to be $\mathbb R^2$ or $S^2$ (a sphere).

\begin{definition}[Embedding]
	A drawing of $G$ on a surface $S=Z_g$ is \textbf{embedding} $G\hookrightarrow S$ if there are no edge crossings.
\end{definition}

\begin{proposition}
	Every finite graph $G$ admit an embedding $G\hookrightarrow Z_g$ for some $g\ge 0$.
\end{proposition}

\begin{proof}
	If $G$ is planar, $G\hookrightarrow Z_0$.

	If $G$ is not planar, for every crossing of edges, we can add a handle to the surface to remove the crossing, so that $g$ increases by $1$.

	Since $G$ is finite, we can repeat this process finitely many times to obtain an embedding $G\hookrightarrow Z_g$ for some $g\ge 0$.
\end{proof}

\begin{definition}[Genus]
	The \textbf{genus} of a graph $G$ is the minimum $g\ge 0$ such that $G\hookrightarrow Z_g$.
\end{definition}

Here we give (without proof) the genus of complete graphs and complete bipartite graphs.

\begin{theorem}[Ringel-Youngs' Theorem]
	$ $

	For $m,n\ge 2$, the genus of $K_{m,n}$ is
	\[
		g(K_{m,n}) = \left\lceil \frac{(m-2)(n-2)}{4} \right\rceil.
	\]

	For $n\ge 3$, the genus of $K_n$ is
	\[
		g(K_n) = \left\lceil \frac{(n-3)(n-4)}{12} \right\rceil.
	\]
\end{theorem}

\section{Euler's Formula}

\begin{theorem}[Euler's Formula]
	For any planar graph drawn on a plane with facet set $F$,
	\[ |V| - |E| + |F| = 2. \]
	Therefore, the number of faces is independent of the drawing of the graph.
\end{theorem}
\begin{proof}
	Assume the graph is connected.
	Fix $|V|$ and we prove by induction on $|E|$.

	Base case: $|E| = |V|-1$. Then the graph is a tree and has $1$ face. So $|V| - |E| + |F| = |V| - (|V|-1) + 1 = 2$.

	Inductive step: Assume the formula holds for $|E| = k$. For $|E| = k+1$, remove an edge $e$ from a cycle, as $G$ is not a tree.

	Applying the inductive hypothesis to the new graph, we have
	\begin{align*}
		|F_G| & = |F_{G\setminus e}| + 1             \\
		      & = (2 - |V| + |E_{G\setminus e}|) + 1 \\
		      & = 2 - |V| + (|E_{G\setminus e}| + 1) \\
		      & = 2 - |V| + |E_G|.
	\end{align*}

	In conclusion, $|V| - |E| + |F| = 2$ holds for all connected planar graphs. For disconnected planar graphs, we can apply the formula to each connected component and sum them up, which still yields $|V| - |E| + |F| = 2$.
\end{proof}


\begin{corollary}
	For a simple planar graph $G$ with $|V|\ge 3$, we have $|E|\le 3|V|-6$.
\end{corollary}
\begin{proof}
	Since $G$ is simple, each face has at least $3$ edges. Each edge can be adjacent to $2$ faces. Therefore, we have
	\[
		2|E| \ge 3|F|.
	\]
	Substituting this into Euler's formula, we get
	\[
		|V| - |E| + \frac{2|E|}{3} \ge 2.
	\]
	Simplifying, we obtain
	\[
		|E| \le 3(|V|-2) = 3|V|-6.
	\]
\end{proof}

\begin{definition}[Girth]
	\term{Girth} of a graph is the size of the smallest cycle in the graph.
\end{definition}

\begin{corollary}
	For a simple planar graph $G$ with $|V| \ge 3$ and girth $\text{gir}$, then
	\[ |E| \le \frac{\text{gir}}{\text{gir}-2}(|V|-2). \]
\end{corollary}
\begin{proof}
	Similar to the previous corollary, since each face has at least $\text{gir}$ edges, we have
	\[
		\text{gir} \times |F| \le 2|E|.
	\]
	Substituting this into Euler's formula, we get the desired inequality.
\end{proof}

\begin{theorem}[Euler's Formula for Surfaces]
	For a graph $G$ embedded on a surface $Z_g$ with facet set $F$ so that each face is an open disc, we have
	\[ |V| - |E| + |F| = 2 - 2g. \]
\end{theorem}
